% Generated draft: requires mathematical and source review.
\documentclass[11pt]{article}
\usepackage[margin=25mm]{geometry}
\usepackage{fontspec}
\setmainfont{Latin Modern Roman}
\usepackage{amsmath,amssymb,mathtools,unicode-math}
\setmathfont{Latin Modern Math}
\usepackage{enumitem,needspace,xurl,hyperref}
\setlength{\parindent}{0pt}
\setlength{\parskip}{5pt}
\setlength{\emergencystretch}{3em}
\newcommand{\N}{\mathbb{N}}
\newcommand{\Z}{\mathbb{Z}}
\newcommand{\Q}{\mathbb{Q}}
\newcommand{\R}{\mathbb{R}}
\newcommand{\C}{\mathbb{C}}
\newcommand{\F}{\mathbb{F}}
\newcommand{\A}{\mathbb{A}}
\newcommand{\PP}{\mathbb{P}}
\newcommand{\E}{\mathbb{E}}
\newcommand{\eps}{\varepsilon}
\newcommand{\dd}{\,\mathrm{d}}
\begin{document}
\section*{UnsolvedMath ID: 701773}
\textbf{Source problem code:} catalogue entry
\subsection{Definitions and mathematical statement}
Let \(G\) be a finite group, and let \(\chi\in\operatorname{Irr}(G)\) be an irreducible complex character: thus \(\chi(g)=\operatorname{tr}(\rho(g))\) for an irreducible representation \(\rho:G\to\operatorname{GL}(V)\) on a finite-dimensional complex vector space \(V\). Its degree is \(\chi(1)=\dim_{\C}V\), where \(1\) is the identity of \(G\). For \(x\in G\), write \(o(x)\) for the least positive integer \(n\) such that \(x^n=1\).

The question is whether, for every finite group \(G\), every \(\chi\in\operatorname{Irr}(G)\), and every \(x\in G\),
\[
\chi(x)\ne 0\quad\Longrightarrow\quad o(x)\mid\frac{|G|}{\chi(1)}.
\]
Equivalently, under the same hypothesis, must there exist a positive integer \(m\) such that
\[
|G|=m\,o(x)\chi(1)?
\]
The supplied record reports that the implication holds when \(G\) is solvable and that, for arbitrary finite \(G\), the same nonvanishing hypothesis implies
\[
\bigl(o(x)\chi(1)\bigr)^4\mid |G|^5.
\]
These two background results are supplied source claims; their individual references and proofs are not supplied, and they have not been independently verified here.
\subsection{Short English statement}
Must every element on which an irreducible complex character is nonzero have order dividing the group order divided by the character degree? The supplied record labels this question open and reports an affirmative answer for solvable groups, together with the stated fourth-power divisibility bound for arbitrary finite groups.
\subsection{Sources}
\begin{itemize}
\item Local UnsolvedMath registry, permanent ID \(701773\), exact source problem code \texttt{catalogue entry}, title \emph{Kourovka Notebook Problem 20.115}; supplied record updated \texttt{2026-10-16T14:00:00Z}. The status \texttt{open} is a supplied registry label, not an independently verified literature assessment.
\item \emph{The Kourovka Notebook (Unsolved Problems in Group Theory)}, Problem 20.115, attributed to T. Wilde, proposed in 2022. The supplied background identifies the 20th issue (2022) and the 21st issue edition (2026), \texttt{arXiv:1401.0300v48}.
\item The supplied collection description names E. I. Khukhro and V. D. Mazurov as editors of the Kourovka Notebook collection covering Issues 1--21 (1965--2026).
\end{itemize}
\end{document}
